\section{The Mesh and Its Cusp}

The committed substrate is the four-dimensional hyperbolic honeycomb $\{3,4,3,4\}$. Its cells are the $24$-cells of the previous section, glued to their neighbours across shared octahedral faces, and the whole assembly tiles negatively curved four-space, growing at its edge. This is the mesh, and the case that it is the one substrate, rather than one choice among many, is made in full later. Here we take it in hand and ask where, inside so strange an object, the ordinary world we inhabit is to be found.

The mesh has two regions that behave very differently, and keeping them apart resolves most of what is puzzling about the model. The interior is the \emph{bulk}, the full four-dimensional hyperbolic space the docks tile. Because hyperbolic space branches like a tree, with each shell holding exponentially more docks than the last, the bulk is vast and tree-like, a place where distances are short through the depth even between points that look far apart on the surface. The boundary is the \emph{cusp}, the flat layer that the curved bulk grows toward at infinity, always one dimension below the bulk. Physical space, the space matter moves in, is the cusp.

\begin{result}[The cusp is three-dimensional space]
At the ideal corners of $\{3,4,3,4\}$ the horosphere, the surface a growing front sweeps out, is flat and three-dimensional, tiled by the cubic honeycomb $\{4,3,4\}$. This flat cubic cusp is ordinary three-dimensional space, and the advance of the growing edge through it is the passage of time. Space is the boundary of the crystal, and time is its growth.
\end{result}

A spinor in the four-dimensional bulk does not vanish on the way to the three-dimensional cusp, it splits. The rotation group of the bulk, $\mathrm{SO}(4)$, restricts to the rotation group of the cusp, $\mathrm{SO}(3)$, and under that restriction a four-component bulk spinor branches cleanly into a pair of two-component spinors, the Pauli spinors of three-dimensional physics. The matter the bulk carries projects onto the cusp as the matter we measure, with its handedness intact. The geometry does the work of dimensional reduction by itself, with no extra rule, because the cusp is built into the honeycomb as its own boundary.

One consequence of this picture is easy to miss and important to state once. The three-dimensional cusp is a thin slice through a four-dimensional space, and a slice has measure zero in the volume it sits in. Almost the entire bulk lies off the cusp, invisible to anything confined to the surface, not by being hidden behind a veil but simply by being in directions a cusp-dweller cannot point. There is therefore genuine room, exactly and without mysticism, for a vast unseen interior, and the later parts of this paper that read the depth of the bulk as the inner life of a mind are reading a region the geometry already provides. What looks like a leap from physics to experience is, structurally, a move from the surface of the crystal into its depth.
